\section{模型的评价与推广}

\subsection{模型的优点}

\begin{itemize}
	\item \textbf{物理口径统一、可复现：}四个问题共享同一套地形净空、等效航程、能耗、充电与通信链路模型，单位、对象编号与计算口径一致，全部计算可由独立 Python 脚本复现，数据与结论可追溯。
	\item \textbf{充分刻画地形与通信约束：}巡航海拔由航段所经过的 DEM 像元最高高程确定，通信状态同时考虑视线遮挡与双向链路预算，较真实地反映了山区环境对运输与通信的双重制约；通信连续性还被写成必须归零的硬指标而非加权项，与题面把「连续通信」列为约束的表述一致。
	\item \textbf{方法选择与问题结构匹配，并以精确解为标尺：}问题一用同质类多重集动态规划在\textbf{允许混用机型}的可行域上求全局最优，并经指派式、集合划分与混用机型三条独立整数规划路径逐行复核（$45/45$、$44/44$，以及 $15$ 区中 $14$ 区一致、S001 因列数超时限如实记为未证最优）；问题二用大邻域自适应搜索（ALNS）配 $\varepsilon$-约束扫描，再用两层小规模精确解标定最优性间隙；问题三用集合覆盖整数规划把悬停高度升为决策变量，并让运输开始时刻进入同一优化问题；问题四用受限增长串完全枚举，枚举结果即全局最优，另以组列集合划分整数规划复核。四问的计算量均在合理时间内完成。
	\item \textbf{揭示多目标权衡：}通过架次数上限—能耗—完成时间的非支配前沿、返航余量—电池能量的二维扰动、通信中断归零的完成时刻代价、以及分区组数—资源规模—均衡度三组定量对照，量化了及时性—能耗—架次、安全性—运力、通信连续性—任务时长、均衡—资源之间的权衡关系，为决策提供了方向性依据。
\end{itemize}

\subsection{模型的缺点}

本文如实列出尚未消除的残留问题，不以「已用高级算法」掩盖其局限。

\begin{itemize}
	\item \textbf{启发式的全局最优性未经证明：}问题二的 ALNS 仍是启发式算法。虽然两层小规模精确验证给出了 Layer A 三个算例间隙全为 $0$、Layer B1/B1$'$ 的最优性证书，但这些算例只含 3 个服务区、8 架无人机对至多 9 个架次，\textbf{机队资源竞争几乎不存在}，故它验证的是组批、能耗、路径与时序逻辑的正确性，\textbf{不能}外推为「全局资源拥塞情形下已证最优」。
	\item \textbf{自由指派层的整数规划间隙很大：}Layer B2 放开了无人机与共享电池的物理指派，求解 90\,s 后最优性间隙仍有 $63.55\%$，只能作为搜索方向上的佐证，\textbf{不构成最优性证明}；同层 B1 的最优解虽已证最优，但该解本身违反 8 项硬时限，故只作下界参考而不作为可行方案。
	\item \textbf{消融实验暴露了一处搜索质量问题：}ALNS 全量配置 ④ 的可行解集严格包含配置 ③，其完成时间（$8041.7$\,s）却比 ③（$7188.7$\,s）劣 $11.9\%$。原因是加权时延归零后目标面变得平坦，而 ④ 的解空间更大，固定迭代预算下更容易停在局部解。本文在 6.4 节如实报告该现象，未将其粉饰为「全量搜索更优」。
	\item \textbf{换站修复产生的新悬停站未经连续精化：}问题三的第三级修复动作为了让某区间「改换一个届时真正在位」的悬停站，会从候选表中直接取站，该站\textbf{未经过} \texttt{refine\_station} 的 $({\rm lon},{\rm lat},h)$ 坐标下降精化。精化会改变站址、进而改变去程时长与建链时刻，因此不能在排班完成后再补做。这些站的有效性由覆盖矩阵的全密度复核保证，但接入裕量不如精化过的站；最终方案一律以 $\Delta t=1$\,s 的逐时刻审计为准。
	\item \textbf{分区方案在给定库存下均不可行，且分区本身在放大资源需求：}问题四的结论是 $K=2$ 与 $K=3$ 在现有库存下都不可行（63 / 301 个分区中库存可行分区数均为 0）。这不是搜索不足，而是题面「各类资源不得跨组调配」这一规则的直接后果——每组必须自备中继护航，使中继无人机需求至少为组数 $K$，而主力组自身峰值就需 2 架。也就是说，\textbf{分区在放大中继与 B 型机队的需求}，本文只能给出补配清单而非可行分区。
	\item \textbf{部分参数为简化假设：}未考虑风场、雨雾、起降排队、充电桩数量等实际因素对飞行与周转的影响。
\end{itemize}

\subsection{模型的改进}

\begin{itemize}
	\item 对问题二可延长迭代预算或采用多起点并行搜索，以缓解上述 ④ 的搜索质量问题；对 Layer B2 可加入割平面或列生成以缩小 $63.55\%$ 的间隙，把自由指派层的最优性也纳入验证范围。
	\item 对问题三可把换站后的新站补做连续精化，但精化必须与排班\textbf{联合迭代}而非事后追加，否则站址变化会使已排定的建链时刻失效。
	\item 对问题四可引入跨组资源共享协议（如资源的时分复用或动态调配许可），以降低「不得跨组调配」带来的机队放大效应；也可在目标中显式加入工作量均衡项，以避免多分组反而更不均衡的现象。
	\item 可建立充电桩数量约束下的多队列充电仿真，细化电池周转模型。
\end{itemize}

\subsection{模型的推广}

本文的“地形感知飞行＋能量余量＋通信链路”联合框架，可推广至其他山区/海岛/灾区的无人机物流、电力巡检、应急通信中继等场景；任务分区与资源配置核算方法也可用于多作战单元、多救援分队的兵力/资源编配决策。
